% Generated draft: requires mathematical and source review.
\documentclass[11pt]{article}
\usepackage[margin=25mm]{geometry}
\usepackage{fontspec}
\setmainfont{Latin Modern Roman}
\usepackage{amsmath,amssymb,mathtools,unicode-math}
\setmathfont{Latin Modern Math}
\usepackage{enumitem,needspace,xurl,hyperref}
\setlength{\parindent}{0pt}
\setlength{\parskip}{5pt}
\setlength{\emergencystretch}{3em}
\newcommand{\N}{\mathbb{N}}
\newcommand{\Z}{\mathbb{Z}}
\newcommand{\Q}{\mathbb{Q}}
\newcommand{\R}{\mathbb{R}}
\newcommand{\C}{\mathbb{C}}
\newcommand{\F}{\mathbb{F}}
\newcommand{\A}{\mathbb{A}}
\newcommand{\PP}{\mathbb{P}}
\newcommand{\E}{\mathbb{E}}
\newcommand{\eps}{\varepsilon}
\newcommand{\dd}{\,\mathrm{d}}
\begin{document}
\section*{UnsolvedMath ID: 30001327}
\textbf{Source problem code:} catalogue entry
\subsection{Definitions and mathematical statement}
Let \(G=(G_t)\) denote the Gaussian process in the source. Centered stationary increments mean that each increment has mean zero and that, whenever the indicated times belong to the time domain,
\[
\mathbb E[G_{t+h}-G_t]=0,
\qquad
(G_{t_j+h}-G_{s_j+h})_{j=1}^{m}
\overset{d}{=}
(G_{t_j}-G_{s_j})_{j=1}^{m}
\]
for every positive integer \(m\), every collection of pairs \((s_j,t_j)\), and every admissible common shift \(h\). This does not impose \(G_0=0\). The time domain and covariance structure are not specified in the supplied record.

To make the term power variation explicit, a possible regular-grid convention for a real-valued process \(Y\) is
\[
V_n^{(p)}(Y;t)
:=\sum_{i=1}^{\lfloor t/\Delta_n\rfloor}
\left|Y_{i\Delta_n}-Y_{(i-1)\Delta_n}\right|^p,
\qquad p>0,\quad t\geq0,\quad \Delta_n\downarrow0.
\]
Here \(p\) is the power, \(\Delta_n\) is the observation spacing, and the process must be defined at the displayed observation times. This is an explicit template, not a sampling convention or range of powers established by the supplied source: neither the meaning of power variation nor the observation regime is fully specified there.

\textbf{Gaussian-process question.} Determine the weak limit of appropriately normalized power variations of \(G\). Under the displayed template, the unspecified target would have the form
\[
a_n\bigl(V_n^{(p)}(G;t)-b_n(t)\bigr)
\xrightarrow{\ d\ }L_p(t),
\]
where the scale \(a_n>0\), any centering \(b_n(t)\), the admissible powers, and the limit law must be specified as part of a precise formulation. The record does not state whether convergence is sought at a fixed time, jointly at finitely many times, or in a function space. No particular normalization or existence of such a limit is asserted here.

\textbf{Volatility-modulated moving-average question.} For the process formally given by
\[
X_t=\int_{-\infty}^{t}g(t-s)\sigma_s\,\mathrm dW_s,
\]
where \(g\) is the kernel, \(\sigma\) is the volatility process, and \(W\) is the driving noise, determine the asymptotic behavior of its power variations and establish a central limit theorem under the intended hypotheses. The integral includes the entire past. Under a usual It\^o interpretation it requires a Brownian driver defined over that past, a suitable filtration and predictable integrand, and conditions ensuring existence of the improper stochastic integral. These are interpretive requirements, not hypotheses supplied by the record. In particular, the source gives no regularity, integrability, dependence, or nondegeneracy conditions on \(g\) and \(\sigma\), and specifies no independence between \(\sigma\) and \(W\).

The central limit theorem target requires a specified normalization, centering, limiting variance or covariance, and mode of convergence; a Gaussian limit conditional on volatility and an ordinary Gaussian limit are distinct possibilities that the supplied formulation does not resolve. Consequently, the record defines a research question but does not determine a single fully quantified limit statement.
\subsection{Short English statement}
Determine the weak limits of normalized power variations for Gaussian processes with centered stationary increments, and determine power-variation asymptotics and a central limit theorem for the displayed volatility-modulated moving average. The supplied record leaves the sampling regime, powers, normalization, process hypotheses, and convergence mode unspecified, so a definitive theorem statement requires further source clarification.
\subsection{Sources}
\begin{itemize}
\item Supplied local UnsolvedMath registry record, permanent ID \(30001327\), title: \emph{Power Variations of Volatility-Modulated Gaussian Processes}; exact source problem code: \texttt{catalogue entry}. Collection: Oberwolfach Reports Open Problems. No individual proposer is identified.
\item Supplied source citation: \emph{Challenges in Statistical Theory: Complex Data Structures and Algorithmic Optimization} (2010), Oberwolfach Reports, pp.\ 2179--2234, DOI \texttt{10.4171/owr/2009/39}. Source URL and supplied literature entry labelled ``original Oberwolfach report 4082'': \url{https://doi.org/10.4171/owr/2009/39}. The registry records proposed year 2009; the citation gives 2010. Their relationship is not resolved by the supplied record.
\item Supplied literature entry labelled ``power variation theorem'': \url{https://www.sciencedirect.com/science/article/pii/S0304414908001427}. No authors or detailed bibliographic citation are supplied.
\item Supplied literature entry labelled ``functional Breuer--Major theorem'': \url{https://arxiv.org/abs/1808.02378}.
\item Literature assessment dated exactly \texttt{2026-08-21}: the supplied curation labels the problem partially solved and states that the fixed-chaos Malliavin criterion and several Gaussian power-variation central limit theorems are established, while the general stationary-increment and stochastic-volatility questions in the extracted passage are not fully classified. This is supplied dated curation, not an independently verified literature assessment; the source requests independent verification of time-sensitive status and cited claims.
\item The registry labels the statement \texttt{corrected\_verified} and reports that its mathematical content was checked against the official Oberwolfach report, with notation, context, and question form normalized. This is a supplied verification claim; no source visit or independent verification is asserted in this draft.
\end{itemize}
\end{document}
